\documentclass{article}
\usepackage{../header}
\usepackage{datetime2}
\usepackage{subcaption}
\def\Xint#1{\mathchoice
	{\XXint\displaystyle\textstyle{#1}}%
	{\XXint\textstyle\scriptstyle{#1}}%
	{\XXint\scriptstyle\scriptscriptstyle{#1}}%
	{\XXint\scriptscriptstyle\scriptscriptstyle{#1}}%
	\!\int}
\def\XXint#1#2#3{{\setbox0=\hbox{$#1{#2#3}{\int}$ }
		\vcenter{\hbox{$#2#3$ }}\kern-0.6\wd0}}
\def\ddashint{\Xint=}
\def\dashint{\Xint-}
\usepackage[
colorinlistoftodos,
textsize=small,          % Sets the text to the smallest LaTeX font size
textwidth=0.7in        % Sets note width slightly smaller than the 1-inch margin
]{todonotes}
\setlength{\droptitle}{-2cm}
\rhead{Summary Nonlinear Wave Equations}
\om 
\title{A Summary of Nonlinear Wave Equations}
\author{\me}
\date{Fall 2026}

\begin{document}
	\maketitle
	\fullline 
	\tableofcontents
	\xappto\thefootnote{\relax}
	\footnotetext{Last updated: \today\ at \DTMcurrenttime}
	\halfline 
	\newpage 
	\setlength{\parskip}{0.55em}
	\section{Overview}
	This document is meant to be a very short and concise summary of the main ideas, results, and tecniques developed in the main notes. The emphasis is not on including every material/result studied, but rather the general ideas and intuition (thus representing the ``heuristics'' behind the theory of nonlinear wave equations and its applications). This document is meant to be read together with the actual set of notes. 
	
	\newpage 
	\part{Nonlinear Wave Equations} 
	\section{Energy Estimates} 
\end{document}